\documentclass[10pt]{article}
\usepackage[margin=20mm]{geometry}
\usepackage{amsmath,amssymb,mathtools,fontspec,unicode-math,xurl,hyperref}
\IfFontExistsTF{DejaVu Serif}{\setmainfont{DejaVu Serif}}{\setmainfont{Latin Modern Roman}}
\setmathfont{latinmodern-math.otf}
\setlength{\parindent}{0pt}\setlength{\parskip}{6pt}\sloppy
\title{Hypothesis Gaps and Core Proof Dependencies
in Intersection Homotopy Theory}
\author{AI-assisted technical review}
\date{}
\begin{document}
\maketitle
Abstract: This report examines arXiv:2211.06096v2 of Chataur, Saralegi-Aranguren and Tanré. Two explicit examples disprove the stated generality of the Van Kampen theorem and of fundamental-group invariance under coarsening. The failures concern ordinary versus perverse connectivity and unrestricted codimensional perversities, respectively. Each is assessed as Major for the named theorem. Neither example refutes the GM-perversity main theorem or the normal Thom–Mather application. The coefficient-domain defect on a cone is Minor because its inspected application already uses a system on the larger space. Two further Major proof defects concern the horn construction establishing the Kan property and the local-to-global step of the main Hurewicz argument. Explicit repairs support the corrected conclusions; they do not validate the original derivations. Scientific adjudication remains pending.\par
\section*{Identity, version and scope}
The paper is David Chataur, Martintxo Saralegi-Aranguren and Daniel Tanré, Relation between intersection homology and homotopy groups, doi: 10.4171/JEMS/1625 (https://doi.org/10.4171/JEMS/1625). The exact reviewed source is the 28-page author manuscript arXiv:2211.06096v2, submitted on 20 February 2025 at 15:16:18 UTC [source]. Its first page separately says 21 February 2025. The publisher records online publication on 11 April 2025 [record]. The acquired author manuscript is not silently identified with the inaccessible publisher Version of Record. All locations below use this manuscript’s PDF page numbers, which agree with its printed numbers.\par
All source text, including references, was read. The checks cover the allowability definitions, simplicial construction, cone formula, local coefficients, Sections 4–5 and the comparison-map proof of the Hurewicz theorem in Section 6. Source-image inspection is individually attributed in the evidence records. External general-position, local-to-global, fibration and classical homotopy theorems were not all independently reproved. The small-resolution statement on page 4 is explicitly a conjecture, not a proved result being rejected here.\par
Equation reference: allow. Write \(Dp(S)=\operatorname{codim}S-2-p(S)\). The basic condition is \[
 \dim\sigma^{-1}(S)\leq \dim\Delta-2-Dp(S).\] A simplex belongs to \(\mathcal G_pX\) only when it and every face are allowable [1, Definitions 1.12 and 2.1, pp. 7–8]. In particular, if \(Dp=1\) on every singular stratum, the entire two-skeleton is that of \(\operatorname{Sing}(X\setminus\Sigma)\). Consequently \(\pi_1^p(X)=\pi_1(X\setminus\Sigma)\), with compatible basepoints. This elementary observation supplies the first two counterexamples.\par
\section*{Van Kampen requires connectivity of allowable paths}
Problem 1.  Theorem 4.1, p. 14, assumes that the two open sets and their intersection are ordinarily path-connected. These assumptions do not imply its surjectivity assertion for perverse fundamental groups.\par
\subsection*{A counterexample satisfying the stated assumptions}
Take the underlying space \(X=\mathbb R^2\) and \[S=\{(0,y):-1\leq y\leq1\},\qquad
 X_0=\{(0,-1),(0,1)\},\quad X_1=S,\quad X_2=X.\] The open segment is a one-dimensional stratum, the endpoints are zero-dimensional strata, and \(X\setminus S\) is regular. This is even a CS set: the middle of the segment has link \(S^0\), and each endpoint has link a circle with one marked singular point. Normality is not required by Theorem 4.1 and does not hold here.\par
Set \(p=-2\) on the open segment and \(p=-1\) at both endpoints, and \(p=0\) on the regular stratum. Then \(p\leq t\) and \(Dp=1\) on all singular strata. No infinite perversity value is needed. Use the open cover \[U_0=\mathbb R^2\setminus\{(0,y):y\geq\tfrac12\},\qquad
 U_1=\mathbb R^2\setminus\{(0,y):y\leq-\tfrac12\}\] and basepoint \(x_0=(1,0)\). The removed rays are disjoint, so the sets cover \(X\). Both sets are path-connected. Their intersection is path-connected as well: its two sides are joined through the gap \(\{(0,y):-\tfrac12<y<\tfrac12\}\). These are exactly the ordinary connectivity hypotheses on p. 14.\par
The regular parts are different: \[U_0\setminus S=\mathbb R^2\setminus\{(0,y):y\geq-1\},\qquad
 U_1\setminus S=\mathbb R^2\setminus\{(0,y):y\leq1\}.\] Each is a slit plane and is simply connected. In contrast, \(\mathbb R^2\setminus S\) deformation retracts onto the circle of radius \(2\). An explicit homotopy is \[H_t(z)=\left(1-t+\frac{2t}{\lVert z\rVert}\right)z .\] For points on the vertical axis in the domain, the radius is strictly greater than \(1\) throughout this homotopy; off that axis the homotopy cannot meet \(S\). Thus it is well-defined and fixes the target circle.\par
Using [allow] (\#allow) through dimension two gives \[\pi_1^p(U_0,x_0)=\pi_1^p(U_1,x_0)=1,\qquad
 \pi_1^p(X,x_0)=\mathbb Z .\] The asserted map is therefore \(1*1\longrightarrow\mathbb Z\), which is not surjective. The missing feature is visible directly: \((U_0\cap U_1)\setminus S\) has two components.\par
\subsection*{Impact, counterevidence and repair}
The proof on p. 14 joins subdivision vertices to the basepoint and says these paths can be chosen in the regular part. Ordinary path-connectivity does not supply those paths. The failure affects the actual conclusion of the stated Van Kampen theorem, a substantive basic result advertised in the paper, so the assessment is Major for that generality.\par
A supported sufficient correction is to require that the regular parts of \(U_0,U_1,U_0\cap U_1\) are path-connected, together with the finiteness needed for the stated subdivision argument (satisfied for CS sets). Then the connecting paths can be chosen allowable, and the usual subdivision proof applies. A groupoid formulation may admit weaker hypotheses but is not proved in this report. As counterevidence to a broader negative inference, Corollary 4.2 assumes normal Thom–Mather structure; its connected open sets have connected regular parts by the paper’s Remark 1.6. This example does not refute that application.\par
\section*{Codimensional perversities need compatibility with coarsening}
Problem 2.  Theorem 4.6, p. 17, asserts fundamental-group invariance for any codimensional perversity \(p\leq t\) when there is no exceptional stratum. This range is false; the displayed induced map need not even be defined by the stated coarsening.\par
\subsection*{An explicit refinement of an axis}
Let \(A=\{(0,0,z):z\in\mathbb R\}\) in \(\mathbb R^3\). Give \(Y\) the singular stratum \(A\) of codimension two. Give \(X\) the finer stratification whose singular strata are \(A\setminus\{0\}\) (two rays) of codimension two and \(\{0\}\) of codimension three. Explicitly, \(Y_0=\varnothing\), \(Y_1=Y_2=A\), whereas \(X_0=\{0\}\), \(X_1=X_2=A\); both top skeleta are \(\mathbb R^3\). These are CS sets without codimension-one strata. At the refined origin, the compact link is \(S^2\) with its two poles as singular points. The identity \(\nu:X\to Y\) is a coarsening, the regular parts agree, and no singular stratum becomes regular.\par
Choose the same codimensional function on both: \[p(2)=-1,\qquad p(3)=1.\] It is zero on regular strata and can be extended, for example, by \(p(1)=-1\) and \(p(c)=c-2\) for \(c\geq4\). It satisfies \(p\leq t\). Its complementary values are \[Dp(2)=1,\qquad Dp(3)=0.\] For \(Y\), all one- and two-simplices avoid \(A\), whence \[\pi_1^p(Y)=\pi_1(\mathbb R^3\setminus A)=\mathbb Z.\] For \(X\), all one-simplices still avoid \(A\), so the regular-part group surjects onto \(\pi_1^p(X)\). Its meridian generator bounds a transverse disk in the plane \(z=0\). Triangulate that disk so that the origin is in the interior of one triangle and on no edge. The resulting two-simplices avoid the open rays. The unique nonempty inverse image of the origin is a point in a two-simplex, which satisfies \(0\leq2-2-Dp(3)=0\); every face avoids the singular set. The disk is an allowable null-homotopy. Therefore \[\pi_1^p(X)=1\ne\mathbb Z=\pi_1^p(Y).\]\par
The same disk simplex also shows that \(\nu\) does not send every \(p\)-full simplex of \(X\) to a \(p\)-full simplex of \(Y\). In the latter stratification the upper bound at the axis is \(2-2-Dp(2)=-1\). Equivalently, Proposition 2.5’s condition \(\nu^*Dp\leq Dp\) fails at the origin: \(1\nleq0\). Thus the problem is not merely a noncanonical identification of two isomorphic groups.\par
\subsection*{Impact and supported restriction}
This is a false invariance assertion under all the printed hypotheses, not a missing explanation of a true isomorphism. It is assessed as Major for Theorem 4.6’s general codimensional range. The counterexample is not a GM perversity: \(p(2)\) is negative and \(Dp\) decreases. In particular, it does not refute Theorem A or Theorem 5.1, both of which explicitly impose the GM condition.\par
A supported repair is to restrict Theorem 4.6 to GM perversities (and the accompanying absence of codimension-one strata). Their complementary perversity is nondecreasing; since coarsening decreases codimension, \(\nu^*Dp\leq Dp\) follows. The source-stratum argument can then establish injectivity, while common regular paths give surjectivity. More general compatible perversities may be possible, but the strongest such range is not established here. The reference to GM perversities immediately after the theorem is counterevidence against treating its intended GM application as disproved; it does not impose a GM hypothesis on the preceding broader statement.\par
\section*{Coefficient domain and core proof dependencies}
Problem 3 (Coefficient system on the wrong member of a pair).  Proposition 2.10, p. 10, quantifies over every local system \(E\) on \(\mathcal G_pL\) while forming \(H_j(\mathcal G_p\mathring cL,\mathcal G_pL;E)\). An arbitrary such system need not extend to the first space.\par
Take \(L=S^1\) with trivial filtration and \(p(v)=0\) on the apex of its two-dimensional cone, so \(Dp(v)=0\). On \(\mathcal G_pL=\operatorname{Sing}S^1\) choose the rank-one integral local system with monodromy \(-1\). The meridian is null-homotopic in the Gajer cone by the same interior-point disk argument used in the cone formula. Its fundamental group is trivial. Every local system on this connected simply connected simplicial set has trivial monodromy when restricted to \(S^1\). The chosen sign system cannot extend, and the displayed relative group is not defined with those data.\par
The supported correction is to take \(E\) on \(\mathcal G_p\mathring cL\) and use its restriction on \(\mathcal G_pL\). The common low-dimensional skeleta then give the stated vanishing range. The application in Proposition 5.2, pp. 19–20, already starts with a coefficient system on the larger space, so this defect does not invalidate that use. The assessment is Minor, a real coefficient-domain restriction with the inspected application preserved.\par
Problem 4 (The specified horn projection).  In the proof of Proposition 2.3, p. 8, radial projection from the barycentre of the missing face is asserted to be an affine map from the whole simplex to the horn. That specified projection is undefined at its centre and has no continuous extension there in dimension two.\par
For a triangle, approach the midpoint of the missing edge along rays aimed at two distinct points of the other two edges. The projection has two different constant limits. This uses a trivially stratified target, so no perversity condition removes the difficulty.\par
There is a direct repair. In barycentric coordinates \((t_0,\ldots,t_\ell)\), with face \(k\) missing, put \(m=\min_{i\ne k}t_i\) and \[\rho_i(t)=t_i-m\ (i\ne k),\qquad
 \rho_k(t)=t_k+\ell m .\] This is a continuous piecewise linear retraction onto the horn, fixing it pointwise. On the missing face it is a piecewise linear homeomorphism onto the horn: if \(y\) is in the horn, the inverse has \(t_k=0\), \(t_i=y_i+y_k/\ell\) for \(i\ne k\). On each linear piece the ambient fibres have dimension at most one; on the missing face they have dimension zero. These facts give precisely the allowability bounds for the whole simplex and each face. The Kan conclusion survives with this replacement. The original construction is nevertheless assessed as Major for its foundational proof dependency. Page 8 calls Proposition 2.3 crucial for the definition of the intersection homotopy groups. Page 12 uses the Kan property to represent homotopy classes and their homotopies by individual full simplices; Proposition 3.5 and the subsequent homotopy arguments use this representation. Page 2 explains why the paper supplies proofs of these basic properties in its general topological setting. The specified discontinuous map cannot supply the required horn fillings.\par
This finding invalidates that derivation, not the Kan conclusion: the explicit PL construction above restores the necessary allowability bounds. Repair success is distinct from severity. No counterexample to the existence of intersection homotopy groups is asserted.\par
Problem 5 (The purported local basis omits angular neighborhoods).  The first step in the proof of Theorem 6.4, p. 26, calls the sets \(C\times\mathring c_tL\) and \(C\times L\times(a,b)\) a basis for a conical chart. They are not a basis in general.\par
Take \(k=2\), \(L=S^2\) with trivial filtration, and the chart \(\mathring cS^2\) with its apex as singular stratum and \(p(v)=0\). The link is simply connected, so the theorem’s link hypotheses hold. An open set \(O\times(a,b)\) away from the apex, with \(\varnothing\ne O\subsetneq S^2\), contains no member of the stated family: each of those sets contains every angular direction at its radii. Thus the assertion that their unions give the full topology, and the subsequent appeal to arbitrary \(W\cap V\), are not justified as written.\par
The existing depth induction supplies a repair. In a chart with central stratum \(S\), use all open subsets of \(V\setminus S\), whose depth is smaller, together with product tip neighborhoods \(C\times\mathring c_tL\). Compactness of \(L\) makes these a basis at the tip, and the former sets cover the off-tip topology. Finite intersections either remain tip neighborhoods or lie off the tip, so the inductive hypothesis and Proposition 6.6 supply \(P_k\) on each. Proposition 6.7 gives finite unions; compactness of finite chains gives the directed-union passage. The Zorn argument can then be used. Before this repair, however, the first step establishes \(P_k\) only for the indicated angularly saturated sets. The second step requires \(P_k(W\cap V)\) for an arbitrary maximal open set \(W\) and a chart \(V\), and this intersection need not belong to that class. The false basis claim therefore leaves Theorem 6.4 unproved by the printed argument. On p. 26, Theorem 6.4 is then used explicitly to deduce Theorem 6.1(2), stated as the main Theorem B(ii) on p. 3.\par
The assessment is Major for this invalid derivation on the main Hurewicz proof path. The depth-induction repair supports the corrected result without changing that assessment. The example refutes the basis claim; it does not refute the Hurewicz conclusion.\par
\section*{Other checked dependencies and limits}
The cone formula preserves homotopy groups through \(Dp(v)\) and kills higher groups; it is not an assertion that the Gajer-space homology always equals intersection homology. The examples involving a torus and \(\mathring c(S^2\times S^3)\) correctly distinguish those theories. For Proposition 6.6, when the cutoff \(q=Dp(v)\) is below \(k\), the assumed connectivity of the link makes the Gajer cone weakly contractible. When the required degree is below the cutoff, the comparison is inherited from the link. This accounts for both the isomorphism and epimorphism ranges, rather than silently extending an isomorphism one degree too far.\par
In Example 6.8, pp. 26–27, the links \(\mathbb{RP}^3,\mathbb{RP}^7\) have fundamental group \(\mathbb Z/2\), so the missing link hypothesis is genuinely absent. The map \(H_3(\mathbb{RP}^7)\to H_3(\mathbb{RP}^{\infty})\) is an isomorphism. The displayed Mayer–Vietoris diagram gives \(J_3\) an isomorphism and \(\operatorname{coker}J_4\cong\mathbb Z/2\). This supports the claimed necessity of the link condition. The double-suspension example concerns a change of regular part and does not contradict the restricted topological-invariance theorem.\par
The generality of ancillary subdivision statements beyond locally finite CS settings has not been exhaustively certified. No additional false theorem is inferred from that limit. A bounded search of the title, publisher record and arXiv record identified no separate correction; a related later coarsening paper was located but is not treated as a correction or as proof of author adoption. The acquisition manifest’s supplement/code searches and the current record checks found no separate paper-specific package. Absence of code is not a defect in this paper.\par
\section*{Conclusion and evidence}
Item | Assessment | Exact effect\par
Problem 1 | Major | Van Kampen surjectivity fails under the printed connectivity assumptions.\par
Problem 2 | Major | Broad codimensional invariance fails; the claimed simplicial map may not exist.\par
Problem 3 | Minor | Arbitrary link coefficients need not extend to the cone.\par
Problem 4 | Major | Invalid horn construction in the Kan proof; an explicit PL replacement restores it.\par
Problem 5 | Major | False basis claim leaves the main Hurewicz local-to-global derivation incomplete.\par
The report-level assessment is Major for the two false general theorem assertions and the two invalid derivations on the Kan and main Hurewicz proof paths. The coefficient-domain defect in Problem 3 remains Minor because the inspected application already uses the correct ambient coefficient system. Problems 1 and 2 do not refute the GM main theorem or the normal Thom–Mather application; Problems 4 and 5 establish proof defects, not counterexamples to the Kan or Hurewicz conclusions. The specified repairs are analytically supported; their availability does not lower the severity of the affected core derivations. Author adoption is unknown. Scientific decision remains pending; the current independent acceptance status is recorded in the verification metadata.\par
\section*{Supporting checks}
The main conversation performed the overall scientific assessment and wrote this report. Their bounded notes, real roles and image hashes are retained. These are AI checks, not human reviews or the external acceptance stage. No independently measured backend equivalence is asserted. Human-review records are empty.\par
The evidence package binds the exact source and manifest, definitions, counterexamples, repair proofs, coverage ledger, current TeX and actual compiler output. Every final report page must be actually viewed before delivery is marked complete. Historical drafts, render-only false states and failures remain preserved separately. No author code was run and no empirical computation is presented as a proof; the counterexamples above are analytic.\par
9 D. Chataur, M. Saralegi-Aranguren and D. Tanré, Relation between intersection homology and homotopy groups, arXiv:2211.06096v2, 20 February 2025, exact acquired author manuscript. https://arxiv.org/abs/2211.06096v2. EMS Press, publisher metadata record for the preceding paper, doi:10.4171/JEMS/1625, online publication 11 April 2025. https://ems.press/journals/jems/articles/14298708.\par
\end{document}
